\documentclass[11pt,a4paper]{article}

\usepackage[utf8]{inputenc}
\usepackage[T1]{fontenc}
\usepackage{lmodern}
\usepackage[margin=1in]{geometry}
\usepackage{amsmath,amssymb}
\usepackage{microtype}
\usepackage{enumitem}
\usepackage{hyperref}
\usepackage{xcolor}
\usepackage{booktabs}
\usepackage{tabularx}

\hypersetup{
  colorlinks=true,
  linkcolor=black,
  citecolor=black,
  urlcolor=blue!50!black,
  pdftitle={A Mathematical-Consistency Verifier for a Phenomenological Curved-Spacetime Quantum-Channel Model},
  pdfauthor={Kwansub Yun}
}

\title{%
  A Mathematical-Consistency Verifier\\
  for a Phenomenological Curved-Spacetime\\
  Quantum-Channel Model\\[0.5em]
  \large Executable Internal-Consistency Checks for the QSOT2 Line
}

\author{Kwansub Yun \\ \small Flamehaven Initiative}
\date{2026-06-19 \\ \small Source: \url{https://github.com/Flamehaven-Labs/QSOT2-Compiler}}

\begin{document}
\maketitle

\begin{abstract}
We present \texttt{QSOT2}, a bounded \emph{mathematical-consistency verifier} for a
phenomenological model mapping classical curved-spacetime curvature norms (Frobenius
norms of the Riemann and Ricci tensors) to single-qubit decoherence channels. The
mapping is a conceptual ansatz with a free calibration parameter $\alpha$; QSOT2 makes
no first-principles or external-physical-validity claim. Its sole deliverable is to
verify that the model holds together as quantum mechanics and as executable software:
temporal-state axioms (linearity, CPTP completeness, trace preservation,
conditionability, density validity) are checked to machine epsilon, channels are CPTP by
construction, a flat-relative Kirkwood-Dirac comparison is reported as a bounded
diagnostic, and a TTM-inspired trace-distance proxy is split into a detector self-test
and the model trajectory. These are executable regression and internal-consistency
tests for the implemented ansatz, not empirical evidence that the curvature-to-channel
mapping is physically correct; the curvature norms themselves are non-derived toy
proxies, so the product $\alpha\lVert\mathrm{Riemann}\rVert_F$ is a pure number carrying
no physical units. The reference run over six spacetime backgrounds emits a
machine-readable result with a four-field claim boundary and an overall
\texttt{DEGRADED\_PASS} verdict across 35 deterministic checks --- degraded for exactly
one documented reason: the flat-baseline KD optimizer does not converge within its fixed
step budget.

\noindent\textbf{Reproducibility.} One command runs the experiment and one runs the test
suite (50 tests, $\sim$96\% coverage). The retained numerical outputs reproduce the
QSOT-Harness r1 math engine; the exact commands and parity targets are given in
\S\ref{sec:repro}.
\end{abstract}

\tableofcontents
\clearpage

\section{Scope and Background}
\label{sec:scope}
\textbf{QSOT} abbreviates \emph{Quantum State Over Time}. \texttt{QSOT2} is a
mathematical-consistency verifier for a phenomenological curved-spacetime quantum-channel
model: it maps the Frobenius norm of a background's Riemann tensor to a single-qubit
decoherence channel and checks that the resulting quantum-mechanical machinery is
internally consistent, deterministic, and reproducible.

\paragraph{Lineage.} The QSOT line passed through three forms. \emph{QSOT~v1} was an
intentional high-formality slop artifact, formal-looking but effectively non-running.
\emph{QSOT-Harness} was an honest reconstruction that made the model run and added an
epistemic-governance shell; it remains the published combined snapshot. \texttt{QSOT2}
re-scopes that work to its mathematical core --- the model and its verification harness
--- with the governance machinery deliberately lifted out so the verification is legible
on its own terms.

\paragraph{Retained scope.} QSOT2 verifies (i) internal quantum-mechanical consistency of
the implemented maps, (ii) deterministic execution of an end-to-end pipeline, and (iii)
reproducible numerical outputs. What it deliberately does \emph{not} claim is centralized
in \S\ref{sec:limits}.

\paragraph{Claim boundary.} A passing QSOT2 check certifies model-internal mathematical
consistency under stated assumptions \emph{only}. The curvature-to-channel mapping is a
phenomenological ansatz whose curvature inputs are non-derived toy values; the
admissibility-gate and verdict labels are implementation-defined classifications, not
physical observables. Accordingly the checks read as executable regression and
consistency tests, not as empirical support for the ansatz. The full non-claim list is
centralized in \S\ref{sec:limits}.


\section{Phase 0: Temporal Axiom Contract Checks}
\label{sec:phase0}
This section should be the formal core of the paper.

It should state the five temporal-axiom checks explicitly:

\begin{itemize}[leftmargin=1.5em]
\item linearity
\item CPTP completeness
\item trace preservation
\item sequential conditionability
\item density-matrix validity
\end{itemize}

For each check, give the mathematical statement and the decision rule. This is
where QSOT2 earns its identity as a verifier rather than as a narrative
artifact. Numerical deviations should be reported as machine-precision evidence,
not as prose claims.


\section{Phases 1--3: Baselines, Curvature Noise, and Boosts}
\label{sec:phase1to3}
This section should cover the retained model pipeline outside KD and the
memory proxy:

\begin{itemize}[leftmargin=1.5em]
\item flat baselines
\item curvature-noise phases across the retained backgrounds
\item relativistic boost amplification checks
\end{itemize}

The emphasis here should be on internal behavior:

\begin{itemize}[leftmargin=1.5em]
\item purity and entropy evolve consistently with the implemented channels
\item Ricci and Riemann play different roles in the retained design
\item boost terms affect the effective noise parameter in the intended direction
\end{itemize}

This section should not overstate physical meaning. It should explain how the
implemented model behaves, not why nature must behave that way.


\section{Phase 4: Flat-Relative KD Signal}
\label{sec:phase4}
This section should isolate the KD phase and keep its interpretation narrow.

Recommended structure:

\begin{enumerate}[leftmargin=1.5em]
\item define the optimized KD quantity used in the implementation
\item report the flat and de Sitter values separately
\item report the flat-relative difference $\Delta_{\mathrm{KD}}$
\item report optimizer convergence status independently from the value itself
\end{enumerate}

Two rules matter here:

\begin{itemize}[leftmargin=1.5em]
\item raw negativity in an optimized basis is not a curvature proof
\item the reported value must track the best objective, not merely the last step
\end{itemize}

This section should therefore read as a bounded diagnostic report, not as a
contextuality theorem.


\section{Phase 5: TTM-Inspired Memory Proxy}
\label{sec:phase5}
Phase~5 runs a TTM-inspired trace-distance proxy: at each step it accumulates the trace
distance between the actual trajectory and the one-step Markovian channel prediction. It
is a bounded backflow diagnostic, not a transfer-tensor reconstruction. The naming is
deliberately comparative: the full Transfer Tensor Method reconstructs non-Markovian
memory through transfer tensors inferred from process data~\cite{CerrilloCao2014},
whereas QSOT2 uses only a trace-distance proxy in that spirit.

Two trajectories are kept distinct:

\begin{itemize}[leftmargin=1.5em]
\item \emph{Synthetic self-test.} An artificial coherence backflow is injected to confirm
  the detector registers $\mathrm{nm} > 0$ (reference value
  $\mathrm{nm} \approx 1.4\times10^{-4}$ at $\alpha=0.1$). A pass validates the detector
  only.
\item \emph{Model-generated trajectory.} Produced by applying the configured channel
  consistently with no injected revival, this yields $\mathrm{nm} = 0$: the
  phenomenological map is Markovian by construction.
\end{itemize}

The detection threshold and the Ricci norm at which it was evaluated are surfaced with
the result, so the (here non-engaged) curvature scaling is explicit rather than implied.
A zero model-trajectory result is a property of the implemented map, not a deep physical
discovery.


\section{Consolidated Results}
\label{sec:results}
Table~\ref{tab:qsot2_results} consolidates the reference run ($\alpha=0.1$, six
backgrounds). The pipeline runs Phase~0 plus Phases~1--5 (35 deterministic checks) and
returns an overall \texttt{DEGRADED\_PASS}: 34 checks pass and one is degraded (the
flat-baseline KD non-convergence). There are no failures and no governance/audit phases.

\begin{table}[h]
\centering
\small
\begin{tabular}{@{}lcc@{}}
\toprule
Quantity & Value & Note \\
\midrule
Overall verdict & \texttt{DEGRADED\_PASS} & 34 pass / 0 fail / 0 skip / 1 degraded \\
Checks (total) & 35 & Phase 0 + Phases 1--5 \\
$P_{\mathrm{Schw}}/P_{\mathrm{dS}}/P_{\mathrm{AdS}}/P_{\mathrm{EH}}$
  & $0.99943/0.63607/0.67764/0.99887$ & purity \\
$S_{\mathrm{dS}}/S_{\mathrm{AdS}}$ & $0.5501/0.5031$ & von Neumann entropy \\
$\Delta_{\mathrm{KD}}$ & $+0.1230$ & flat-relative, non-converged \\
$\mathrm{nm}$ (self-test / model) & $1.41\times10^{-4}\;/\;0$ & detector vs.\ model \\
axiom deviations & $\le 4.4\times10^{-16}$ & machine epsilon \\
\bottomrule
\end{tabular}
\caption{Consolidated QSOT2 reference-run outputs.}
\label{tab:qsot2_results}
\end{table}

\paragraph{Sensitivity.} The qualitative classification is stable across a manual
three-point $\alpha$ sweep while magnitudes remain calibration-dependent: at
$\alpha\in\{0.01, 0.1, 1.0\}$ the flat-relative signal is
$\Delta_{\mathrm{KD}} = +0.0135,\ +0.1230,\ +0.1235$, and the model-trajectory memory
result stays $0$ throughout.

\paragraph{Parity.} Every retained mathematical quantity above reproduces the
QSOT-Harness r1 math engine numerically. QSOT2 differs only by design: the
governance/audit surfaces are absent and the result schema is leaner.


\section{Limits and Explicit Non-Claims}
\label{sec:limits}
Put the non-claims here, and keep them here.

This section should explicitly state:

\begin{itemize}[leftmargin=1.5em]
\item no first-principles derivation of the curvature-to-channel map
\item no external physical validation
\item no claim that the free calibration parameter is physically derived
\item no claim that the memory proxy is a full transfer-tensor reconstruction
\item no governance-heavy audit interpretation in this paper
\end{itemize}

The rest of the paper should not keep repeating these points unless a local
clarification is necessary. This is the main stylistic lesson from
QSOT-Harness: centralize the non-claims instead of letting them dominate every
section.


\section{Reproducibility}
\label{sec:repro}
This section should be practical and short.

It should contain:

\begin{itemize}[leftmargin=1.5em]
\item the exact command used to run the experiment
\item the exact command used to run the tests
\item the expected artifact filenames
\item the expected schema identifier
\item the retained numeric parity targets relative to the source migration line
\end{itemize}

Avoid environment-dependent ornament such as exact wall-clock timings unless the
timing itself is the subject of the paper.


\section*{Acknowledgements}
Keep this section minimal and accurate.

If AI-assisted coding tools were used for drafting, refactoring, or review
support, say so plainly without implying institutional collaboration. If the
paper depends on open-source scientific libraries, acknowledge those projects
directly rather than broad organization-level attributions.


\end{document}
